\subsection{Finite recognition of the simultaneous ten-point centralizer}

\begin{result}{A.4.34}{Proved}{An explicit finite refinement recognizes the actual Klein-four centralizer at every specified lower quotient}
Retain the actual ten-point marking (A.T10), with the four
points $D_0$ and any initial collection of the fixed pairs
$A_0,B_0,C_0$.  Let $G_t$, for $0\leq t\leq3$, be the
inverse image of the permutation centralizer of the named
$H=C_2^2$ on these $4+2t$ marks; use the order
$C_0,B_0,A_0$ when adding the pairs, so the reverse
forgetting order is $A_0,B_0,C_0$.
Write $N_t=C_{G_t}(H)$.  For every specified finite quotient
$q:G_t\to Q$, the construction below gives a specified
finite quotient $q^+:G_t\to Q^+$ refining $q$, such that
\[
 \operatorname{im}\bigl(C_{Q^+}(q^+(H))\longrightarrow Q\bigr)
                              =q(N_t).                 \tag{HCF.1}
\]
The refinement uses only the actual point-pushing free groups,
the characteristic kernels (BCF.2), finite induced
representations, and joins and quotients of specified finite
images.  In particular it applies to the actual induced
$630$-factor quotient of A.4.32.  It recognizes the image of
the simultaneous centralizer, with both named involutions
retained; it makes no assertion that every finite subgroup
is conjugate to the named $H$ at the input level.
\end{result}

\begin{proof}
\par\medskip\noindent\textbf{The actual ordinary maps.}\quad
Deleting the last added fixed pair gives a surjection
\[
 \pi_t:G_t\longrightarrow G_{t-1},\qquad K_t=\ker\pi_t.
                                                               \tag{HCF.2}
\]
The kernel is the two-point unordered configuration group
of the sphere with its $2t+2$ old punctures.  Its index-two
subgroup $K_t^0$ fixes the two new labels individually.
The inclusion of the actual $H$ is carried isomorphically
by $\pi_t$, since its action on $D_0$ is faithful.
If the new pair is $\{p,q\}$, denote by $a_t\in H$ the
involution with that fixed set.  Thus the three values
of $a_t$ in reverse order are $a,b,c$.

The restriction of (HCF.2) is an actual isomorphism
\[
                         N_t\xrightarrow{\sim}N_{t-1}.
                                                               \tag{HCF.3}
\]
Here the added pair is unordered: the other two elements
of $H$ interchange its members.  To verify (HCF.3), use
the ordinary centralizer/quotient comparison proved in A.1.
An $H$-equivariant sphere with the marked regular quartet
$D_0$ has exactly the three additional fixed pairs, of
which any specified subset may be marked or unmarked.
Adding or deleting the entire fixed set of $a_t$ is inverse
on its equivariant isomorphisms.  In the quotient sphere
the three branch types are still distinguished by the
three named elements of $H$, including when their fixed
points are unmarked upstairs.  The ordinary descended
mapping class is consequently pure on the three branch
points and the quotient of $D_0$.  Conversely every such
pure downstairs class fixes the covering monodromy and
lifts.  Its deck kernel is exactly $H$.
Equivalently, in the explicit family of (A.T13) the
quotient coordinate $u$ and the quartet coordinate satisfy
$\lambda_D=1-1/u$; this is an isomorphism of the parameter
thrice-punctured spheres, independently of which fixed
pairs have been marked.  The deck map is the identity
$H\to H_D$.  Thus forgetting to $D_0$ identifies every
$N_t$ with the same $\Gamma_{0,4}\times H_D$, and
(HCF.3) is the identity on that actual target.
In particular
\[
       G_0=N_0=\Gamma_{0,4}\times H_D,\qquad
                        H_D\subseteq Z(G_0).           \tag{HCF.4}
\]
This supplies the starting case of the finite induction.

\par\medskip\noindent\textbf{A finite obstruction to an isolated transposition.}\quad
For a sphere with $n\geq5$ marks use the actual binary
quotient $E_2(n)$ of A.4.29: its pure kernel is
\[
 V_2(n)=\mathbb F_2^{\binom n2}/
          \langle\operatorname{star}_i,\operatorname{total}\rangle,
                                                               \tag{HCF.5}
\]
and the raw crossing cocycle comes from braid crossings
modulo four.  There is no element of order two in $E_2(n)$
above a single transposition $(p\ q)$.
Indeed choose three other labels $r,s,t$.  In the dual edge
space take the sum of the two triangle-edge vectors
$(p,q,r)$ and $(r,s,t)$.  Denote the resulting functional
by $\ell$.  Every vertex has even degree and the total
number of edges is six.  Thus $\ell$ annihilates every
relation in (HCF.5).  It is invariant under $(p\ q)$,
and it takes the value one on the edge $pq$.
Choose an actual half twist $T_{pq}$ along an arc joining
$p$ and $q$.  Its square is the pure pair twist with class
$e_{pq}$ in (HCF.5).  Write an arbitrary lift of the same
transposition as $vT_{pq}$; its square has pure coordinate
\[
             (1+(p\ q))v+e_{pq};
\]
applying $\ell$ gives one.  The square is therefore
nonidentity, as asserted.  This argument retains the
actual source quotient and its extension cocycle.

\par\medskip\noindent\textbf{The unordered-pair fibre detector.}\quad
We first construct, from the given $q:G_t\to Q$, a finite
quotient $u:G_t\to U$ refining $q$ such that
\[
 \operatorname{im}\bigl(C_U(u(a_t))\cap u(K_t)
                                     \longrightarrow Q\bigr)=1.
                                                               \tag{HCF.6}
\]
Join $q$ with the actual map to $E_2(4+2t)$, and call this
finite image $Q'$.  Let $G_t^0$ be the index-two subgroup
fixing $p,q$ individually.  It contains $a_t$.
Forgetting $p$ gives a surjection
\[
 \pi_p:G_t^0\longrightarrow G_q,
     \quad F_p=\ker\pi_p,
 \qquad
 \pi_q:G_q\longrightarrow G_{t-1},
     \quad F_q=\ker\pi_q.                              \tag{HCF.7}
\]
Both free groups are actual one-point Birman kernels;
their ranks are $2t+2$ and $2t+1$.  The named involution
$a_t$ fixes both base points, so A.4.33 applies twice.

Apply A.4.33 first to $Q'|_{G_t^0}$ and $F_p$, obtaining
$u_1:G_t^0\to U_1$.  Form the literal finite quotient
\[
                 \overline U_1=U_1/u_1(F_p)
                                                               \tag{HCF.8}
\]
of $G_q$.  Apply A.4.33 to this quotient and $F_q$,
obtaining $u_2:G_q\to U_2$.  On $G_t^0$ retain the actual
joint map $(u_1,u_2\pi_p)$, with image $T$.
For $k\in K_t^0$ whose image in $T$ centralizes $a_t$,
its $U_2$ projection lies in $u_2(F_q)$ and centralizes
$a_t$.  The second detector says that its projection
to $\overline U_1$ is trivial.  Consequently $u_1(k)$
lies in $u_1(F_p)$ and centralizes $a_t$.  The first
detector then gives $q'(k)=1$ in $Q'$.
This verifies the whole ordered-pair fibre by the two
specified quotient maps, rather than by a claim about
an abstract intersection of completed groups.

Induce the finite homomorphism $G_t^0\to T$ to $G_t$
using its two right cosets, and join it with $q':G_t\to Q'$.
Let $u:G_t\to U$ be that actual image in
$Q'\times(T\wr\mathfrak S_2)$.  On the subgroup $G_t^0$,
the coordinate at its identity coset is the original map
to $T$.  Thus every $k\in K_t^0$ whose image under $u$
centralizes $a_t$ has trivial image under $q'$.
For $k\in K_t\setminus K_t^0$ with the same commutation
property, its square lies in $K_t^0$ and still centralizes
$a_t$.  It would therefore have trivial image in $E_2(4+2t)$.
But the image of $k$ in that quotient lies above the
single transposition $(p\ q)$, contradicting the preceding
square calculation.  This excludes this second case and
proves (HCF.6) for the unordered pair.
Every element of $u(K_t)$ has an actual representative in
$K_t$, so the same proof applies to the entire displayed
finite centralizer fibre.

\par\medskip\noindent\textbf{The recursive simultaneous comparison.}\quad
Define the refinement $\mathcal R_t(q)$ recursively.
For $t=0$ put $\mathcal R_0(q)=q$, using (HCF.4).
For $t\geq1$ form the just-constructed $u:G_t\to U$ and
its actual quotient
\[
                \overline U=U/u(K_t),                  \tag{HCF.9}
\]
which is a finite quotient of $G_{t-1}$ through $\pi_t$.
Apply the already defined $\mathcal R_{t-1}$ to this
quotient.  Finally let $\mathcal R_t(q)$ be the actual
joint image of $u$ and that refined lower map composed
with $\pi_t$:
\[
 G_t\longrightarrow
 U\times\mathcal R_{t-1}(\overline U),\qquad
 g\longmapsto\bigl(u(g),
             \mathcal R_{t-1}(\overline U)(\pi_t(g))\bigr).
                                                               \tag{HCF.10}
\]
All factors and both projections in this formula are actual
finite quotient maps of the indicated discrete groups.

To prove (HCF.1), represent an element of the centralizer
in (HCF.10) by an actual $g\in G_t$.  The induction
hypothesis supplies $c_0\in N_{t-1}$ whose image in
$\overline U$ agrees with that of $\pi_t(g)$.
The actual isomorphism (HCF.3) supplies $c\in N_t$
with $\pi_t(c)=c_0$.  Thus $u(gc^{-1})\in u(K_t)$,
and it centralizes $u(H)$, hence also $u(a_t)$.
Formula (HCF.6) gives $q(gc^{-1})=1$, so $q(g)=q(c)$.
This proves the inclusion of the left side of (HCF.1)
in the right side.  Every actual $c\in N_t$ centralizes
$H$ at every finite quotient, proving the reverse inclusion.
The induction has therefore identified the exact finite
image of the simultaneous centralizer.

\par\medskip\noindent\textbf{Explicit bounds and finite instructions.}\quad
The construction does not ask an unspecified finite
quotient to separate a desired pair of elements.
For input order at most $m$ and free rank $r$, set
\[
\begin{aligned}
 D(m,r)&=(m!)^{(m!)^r},\\
 W(m,r)&=D(m,r)\,2^{1+D(m,r)(r-1)}
       3^{\,1+D(m,r)2^{1+D(m,r)(r-1)}(r-1)},\\
 B(m,r)&=mW(m,r)^r.
\end{aligned}                                                   \tag{HCF.11}
\]
These are the explicit monotone upper bounds (BCF.4).
The binary join in the pair step has order at most
$m'=m(4+2t)!2^{(4+2t)(1+2t)/2}$.
The subsequent groups satisfy
\[
 |U_1|\leq B(m',2t+2),\quad
 |U_2|\leq B(B(m',2t+2),2t+1),\quad
 |U|\leq2m'(|U_1||U_2|)^2.                              \tag{HCF.12}
\]
The final joint image in (HCF.10) has order at most
$|U|$ times the recursively bounded lower image.
The recursion has only $t\leq3$ steps.
To build each characteristic kernel one enumerates the
finite set of generator images in the prescribed symmetric
group; the free homology covers are then the ordinary
Schreier graphs with their binary or ternary cycle spaces.
Induction to the two point cosets is the usual explicit
Schreier-coordinate map.  Thus (HCF.7)--(HCF.12) specify
finite groups, their source maps, and a finite upper bound
using only finite enumeration and free-word operations.
No computation of every element of these large groups is
asserted here.  The statement is the proved finite
centralizer refinement, not the remaining A.4 pentagon
shadow census.
\end{proof}
